\documentclass[../../main.tex]{subfiles}
\begin{document}
\section{Route S: fixed-eigenfunction localization}
\label{sec:spectral-route}

\paragraph{Thesis.}
Follow a first eigenfunction --- rather than a candidate cut --- through stochastic localization,
and apply the moment-map quadratic control of Section~\ref{sec:family-moment-map} to its whitened
posterior covariance tensor. The target is an absorptive, function-aware source/damping estimate
over a universal amount of localization time.

This route attacks the spectral object that \emph{defines} KLS directly, and it retains the
tensor/covariance orientation that a global operator-norm bound discards --- the property
Section~\ref{subsec:kls-spike-obstruction} showed to be necessary. It shares the localization
backbone of Group~2 with Route~E and differs only in the object it tracks.

The route is \textbf{live, secondary}. Its current entry points and acceptance gates are in the
research control plane's \texttt{routes.md} and \texttt{gating.md}.

\subsection{Setup and the exact fixed-function SDE}
\label{subsec:spectral-sde}

Work first on smooth, strongly log-concave isotropic approximants, where the weighted Laplacian
$L=\Delta-\nabla V\cdot\nabla$ has discrete spectrum. Let $f$ be a normalized first nonconstant
eigenfunction, $-Lf=\lambda f$, and define along the localization of
Section~\ref{subsec:sl-process}
\begin{equation}
  g_t=\Cov_{\mu_t}(f,X),
  \qquad
  H_t=\E_t\bigl[(f-\E_tf)(X-a_t)^{\otimes2}\bigr],
  \qquad
  A_t=\Cov_{\mu_t}(X).
  \label{eq:spectral-quantities}
\end{equation}
Here $g_t\in\R^n$ is the covariance of $f$ with the coordinates, and $H_t$ is the
symmetric-matrix-valued tensor coupling $f$ to the second-order posterior structure.

The exact evolution is
\begin{equation}
  \dd g_t=H_t\dd W_t-A_tg_t\dd t .
  \label{eq:spectral-sde}
\end{equation}
So $\norm{H_t}_{\HS}^2$ is the \emph{source} for $\abs{g_t}^2$ and $2g_t^TA_tg_t$ is its
\emph{exact damping}. This is the fixed-function analogue of the two-color Riccati identity of
Section~\ref{sec:riccati}: source and damping are separated exactly, with no inequality spent.

\subsection{What the July 2026 input gives, and what it leaves}
\label{subsec:spectral-whitened}

Conditional on the version-1 preprint of Theorem~\ref{thm:letwin-qcts}, the eigenfunction tensor
satisfies the optimal \emph{intrinsic} static estimate in whitened coordinates,
\begin{equation}
  \norm{A_t^{-1/2}H_tA_t^{-1/2}}_{\HS}^2\le8\,\Var_{\mu_t}(f).
  \label{eq:whitened-tensor}
\end{equation}
This is the exact analogue, for the fixed-eigenfunction object, of the static quadratic-chaos
input that Section~\ref{sec:qcts} supplies for the fixed-cut object.

Crude unwhitening of \eqref{eq:whitened-tensor} --- multiplying back by $A_t^{1/2}$ on both sides
--- leaves a factor $\lmax(A_t)^2$, and hence exactly the universal-time dynamic alignment
problem that Section~\ref{subsec:sl-where-the-log-lives} identified as the residual cost. So this
route does not escape the difficulty by importing the preprint; it relocates it to an object with
more structure.

\begin{remark}[The extra structure a cut does not have]
\label{rem:spectral-extra-structure}
Unlike an indicator, the eigenfunction satisfies an equation, and that equation fixes two energies
which the cut has no analogue of:
\[
  \E\abs{\nabla f}^2=\lambda,
  \qquad
  \E\norm{\Hess f}_{\HS}^2\le\lambda^2 .
\]
Exploiting these fixed energies inside the posterior channel is the possible advantage of this
route over Route~E. The route succeeds only if it controls the tensor's incidence in inflated
covariance spaces, or consumes the exact damping in \eqref{eq:spectral-sde}. Another global
covariance-norm estimate is not enough, and by Proposition~\ref{prop:covariance-spike} could not
be.
\end{remark}

\begin{lemma}[Posterior eigenfunction-defect calculus]
\label{lem:mm-posterior-defect}
In the regular setting above, realize localization through the planted channel
$c_t=tX+B_t^{\mathrm{obs}}$, where $X\sim\mu$ and $B^{\mathrm{obs}}$ is an independent Brownian
motion, and put
\[
  R_t(x)=(c_t-tx)\cdot\nabla f(x),
  \qquad \rho_t=R_t-\E_tR_t,
  \qquad m_t=\E_tf.
\]
Define
\[
  b_t=\E_t\nabla f,
  \quad C_t=\E_t[(X-a_t)\otimes\nabla f],
  \quad u_t=\E_t[\rho_t(X-a_t)],
  \quad K_t=\E_t[\rho_t(X-a_t)^{\otimes2}].
\]
Then, for every $t\ge0$,
\begin{align}
  \E_tR_t&=\lambda m_t,
  &\lambda g_t&=b_t+u_t,
  &\lambda H_t&=2\operatorname{sym}C_t+K_t,
  \label{eq:mm-posterior-defect-identities}\\
  \E\Var_t(R_t)
  &=t\lambda-\lambda^2\int_0^t\E\abs{g_s}^2\dd s,
  &\E\int_0^t\norm{C_s}_{\HS}^2\dd s
  &\le\lambda-\lambda^2\abs{g_0}^2,
  \label{eq:mm-posterior-defect-budgets}\\
  \E\E_t\abs{\nabla R_t}^2
  &\le t\lambda^2+t^2\lambda.
  \label{eq:mm-posterior-defect-gradient}
\end{align}
These identities do not by themselves unwhiten the high-covariance part of $H_t$.
\end{lemma}

\subsection{The headline open target}
\label{subsec:spectral-headline}

\begin{question}[Fixed-eigenfunction full-damping occupation]
\label{q:mm-spectral-occupation}
On smooth strongly log-concave isotropic approximants, let $-Lf=\lambda f$ be a normalized first
nonconstant eigenfunction and let $g_t,H_t,A_t$ be as in \eqref{eq:spectral-quantities}. Do
universal constants $T_0,C_0,C_1>0$ exist such that, for every $t\le T_0$,
\begin{equation}
  \E\int_0^t\norm{H_s}_{\HS}^2\dd s
  \le C_0t+C_1\E\int_0^t\abs{g_s}^2\dd s
  +\E\int_0^t2g_s^TA_sg_s\dd s,
  \label{eq:spectral-occupation}
\end{equation}
uniformly through approximation?
\end{question}

\begin{question}[Sufficiency bridge]
\label{prop:spectral-sufficiency}
Prove that an affirmative answer to Question~\ref{q:mm-spectral-occupation}, uniform through
regularization, implies a universal spectral gap, hence Conjecture~\ref{conj:kls}.
\end{question}

The intended mechanism is the fixed-function SDE \eqref{eq:spectral-sde} together with
Gr\"onwall, posterior Brascamp--Lieb, and the fixed-gradient density martingale; the shape of the
argument should parallel the two-color analysis in Sections~\ref{sec:mass-martingale}
and~\ref{sec:carleson}.  That mechanism has not yet been written as a proof with constants,
terminal variance control, and a regularization limit, so it is tracked as an open bridge rather
than a proved implication.  The estimate \eqref{eq:spectral-occupation} itself is also open, and
must either hold uniformly through regularization or be reformulated for approximate Rayleigh
minimizers.

The endpoint structure of \eqref{eq:spectral-occupation} is worth naming.  Here the source may be
charged against the \emph{full} exact damping.  In the evolution of $\E|g_t|^2$ those two terms
then cancel, leaving a closed Gr\"onwall inequality from the linear budget and the lower-order
term.  A strict damping surplus would be useful but is not required for the KLS sufficiency
bridge; this differs from the absorption margins required in the two-color route.

\begin{lemma}[Time-weighted fixed-function source budget]
\label{lem:mm-time-weighted-fixed-source}
Let $\mu=e^{-V}\dd x$ be a smooth probability with $\Hess V\succeq\kappa I$, $\kappa\ge0$, and
let $f$ be a fixed square-integrable test for which the posterior identities are justified.  Put
$g_t=\Cov_{\mu_t}(f,X)$, $H_t=\E_t[(f-\E_tf)(X-a_t)^{\otimes2}]$, and $A_t=\Cov_{\mu_t}(X)$.
Then, for every $T>0$,
\begin{align}
 &\E\int_0^T(\kappa+t)\norm{H_t}_{\HS}^2\dd t
 +2\E\int_0^T\left(\abs{g_t}^2-(\kappa+t)g_t^TA_tg_t\right)\dd t \notag\\
 &\hspace{7em}\le \Var_\mu(f)-\kappa\abs{g_0}^2.
 \label{eq:mm-time-weighted-fixed-source}
\end{align}
In particular, for $\kappa=0$,
$\E\int_0^\infty t\norm{H_t}_{\HS}^2\dd t\le\Var_\mu(f)$.
This estimate retains one power of time and by itself gives no unweighted initial-layer bound.
\end{lemma}

\subsection{An equivalent sufficient reformulation}
\label{subsec:spectral-kappa-form}

There is a second, coarser way to state what this route needs. By
Proposition~\ref{prop:letwin-kappa}, $\kappa_n\le2\sqrt2$ conditional on the preprint. Hence
\emph{any} dimension-free comparison of the form
\begin{equation}
  \CP(\mu)\le C\bigl(1+\kappa_n^2\bigr)
  \label{eq:spectral-kappa-sufficient}
\end{equation}
valid for every isotropic log-concave $\mu$ would prove KLS. Equivalently: it suffices to remove
the residual $\sqrt{\log n}$ from the current spectral comparison \eqref{eq:kls-bridge} while
retaining only a universal function of $\kappa_n$. Formulation
\eqref{eq:spectral-kappa-sufficient} is useful as a target statement but supplies no mechanism;
Question~\ref{q:mm-spectral-occupation} is the mechanism-bearing form.

\subsection{The audited \texorpdfstring{$H^{-1}$}{H-1} endpoint}
\label{subsec:spectral-hminus1-endpoint}

An earlier formulation of this route's first lemma was an $H^{-1}$ residual bound. The cycle-1
audit reclassified it, and the reclassification is worth recording because it is exactly the kind
of error the repository's soundness contract exists to catch.

With $b=\int\nabla f\dd\mu$, the proposed target was
\begin{equation}
  R(f)=\sum_i\norm{\partial_if-b_i}_{H^{-1}(\mu)}^2\le C\lambda .
  \label{eq:spectral-residual}
\end{equation}
It is still sufficient for KLS. But it is also quantitatively \emph{implied} by KLS: one has the
two-sided relation
\begin{equation}
  1+\abs b^2-\frac{2\abs b^2}\lambda\ \le\ R(f)\ \le\ 1-\frac{\abs b^2}\lambda .
  \label{eq:spectral-residual-two-sided}
\end{equation}
Its exact heat representation controls the short-time and high-frequency parts at the desired
scale; the long-time low-spectrum tail is the complete KLS-strength residue. Accordingly
\eqref{eq:spectral-residual} is retained as a \emph{calibrated endpoint and diagnostic}, not
advertised as a likely preliminary lemma.

A second cycle-1 finding is recorded with it: a direct unweighting of the positive moment-map
Stein form is \emph{false} on truncated-exponential first eigenfunctions. That is a genuine
refutation of a natural first attempt, and it is why \eqref{eq:whitened-tensor} is stated in
whitened form.

\subsection{Fences and comparison with the other routes}
\label{subsec:spectral-fences}

The obstructions proved elsewhere in Part~III are scoped, and it matters which of them bind here.

\begin{itemize}[leftmargin=1.6em]
  \item The projection ceiling of Section~\ref{sec:qcts} does \emph{not} refute this route: the
  moment map uses information beyond radial and projection tests.
  \item The rank-one product-budget refutation of Section~\ref{sec:product-stress} is specific to
  a fixed-cut localization counterexample and imposes no no-go here.
  \item A universal Lipschitz Gaussian transport is too strong for exponential tails
  (Warning~\ref{rem:no-lipschitz-transport}); the weaker expected-Jacobian criterion
  \eqref{eq:brownian-derivative} is sufficient for KLS but currently reuses KLS-scale input, so
  it ranks below the target of this section.
  \item Classical needles do not preserve the isotropic covariance constraints
  (Section~\ref{sec:family-needles}); that is a structural warning, not a theorem excluding all
  needle arguments.
\end{itemize}

\begin{remark}[Relation to Route~C]
\label{rem:spectral-vs-cmh}
Routes~S and~C both consume moment-map information and both must ultimately handle a
test-dependent object rather than a constant matrix, which is the shared lesson of
Section~\ref{subsec:mm-audit}. They are nonetheless different programs: Route~S keeps stochastic
localization and makes the test dependence dynamic, while Route~C
(Section~\ref{sec:moment-map-cmh}) is deterministic and makes it geometric, through Haar fields on
Schur fibers. Similarity of the residual loss is \emph{not} evidence that the two open statements
are equivalent, and no result merges them.
\end{remark}

\paragraph{Promotion gate.}
The control plane promotes Question~\ref{q:mm-spectral-occupation}, or admits further internal
claims on this route, only when one of the following is established: the absorptive
source/damping estimate \eqref{eq:spectral-occupation} uniformly on regular approximants together
with a passage to arbitrary log-concave measures; or a route-fatal counterexample to every such
function-aware occupation estimate.
\end{document}
